\section{ADR-009: Visual identity (Tape \& Console)}
\textbf{Status:} accepted (M13).

\textbf{Context:} in November 2026 we found StemKit, another MIT desktop app
that turns songs into stems. No code was shared either way, but our dark
interface looked almost the same as theirs: a near-black background with two
radial glows (violet and pink), a 264~px left side bar with a green
\emph{Engine ready} dot at the bottom, a gradient rounded-tile logo with
equaliser bars, violet-to-pink gradients and Segoe UI. Users could mistake one
app for the other. Cover\_Maker also has a clearer audience than a generic
stem tool: musicians who practise and record covers. The identity had to
change before the next release (0.4.0).

\textbf{Decision:} a warm analogue-hardware identity called \emph{Tape \&
Console}, with the descriptor \emph{the cover musician's studio}
(\emph{el estudio del músico de covers}). The name stays Cover\_Maker.
\begin{itemize}
  \item \textbf{Colour.} Warm near-black and cream instead of blue-black and
        white: background \texttt{\#0c0b0a}, surfaces \texttt{\#15130f},
        \texttt{\#1d1a16}, \texttt{\#29251f}, text \texttt{\#f1e9dc}, muted
        \texttt{\#a59d90}, hairline borders at 10~\% of the text colour. One
        accent, signal orange \texttt{\#ff6a3d} (pressed \texttt{\#e04e22}),
        for fills and the active state; blue \texttt{\#5b7cff} only for links;
        status colours \texttt{\#8fd16a}, \texttt{\#ffc04d}, \texttt{\#ff4d4d}.
        No gradients and no background glows. A light ``cream paper'' theme
        uses the same hues; text in the accent uses a darker orange there so
        that every text colour meets WCAG AA in both themes.
  \item \textbf{Stem colours.} Vocals \texttt{\#ffcf5c}, guitar
        \texttt{\#ff8f4d}, drums \texttt{\#5fd3b5}, bass \texttt{\#6fa8dc},
        piano \texttt{\#c8d46a}, other \texttt{\#9b948a}. In the light theme
        stem-coloured text is mixed with the text colour to stay readable.
  \item \textbf{Type.} Archivo SemiExpanded for headings and console labels,
        IBM Plex Sans for text, IBM Plex Mono for readouts. All three are SIL
        Open Font License 1.1 and bundled as WOFF2 files, never loaded from
        the network (the app works offline and makes no third-party requests).
  \item \textbf{Layout.} A 64~px icon rail replaces the side bar (names show on
        hover and keyboard focus). A transport dock at the bottom holds the
        player, so a song keeps playing on every page. Engine and runtime status
        are LEDs in an amber-on-black LCD panel in the dock, next to the clock
        and the slots for tempo, key and tuning that the practice tools will
        fill (M14).
  \item \textbf{Signature parts.} Segmented LED meters with peak-programme
        ballistics (instant attack, 20~dB in 1.7~s fall-back, IEC~60268-10
        type~I, and a held peak), knobs for pan, console faders, buttons with
        an LED that lights when they are on.
  \item \textbf{Logo.} A tape reel whose flange opens into a C; the tape
        leaving through the gap is the underscore of the name. The app icon
        puts it on a round dark chassis, not on a rounded tile. Two other
        directions (a fader cap, and \texttt{C\_} as a blinking cursor) were
        drawn and rejected by the maintainer.
  \item \textbf{Motion.} UI transitions take 150~ms; page content rises
        8--12~px into place with \texttt{cubic-bezier(.22,1,.36,1)} and a
        70~ms stagger; meters are drawn from their ballistics, never eased.
        \texttt{prefers-reduced-motion} turns all of it off.
\end{itemize}

\textbf{Consequences:} the app no longer resembles StemKit, and its look
says what it is for. The bundled fonts add about 0.6~MB and their licences are
listed in \texttt{THIRD\_PARTY\_NOTICES.md}. The persistent player moved from
the mixer page into a shared store (\texttt{lib/transport}), which also made
per-stem pan and keyboard shortcuts possible. The documentation, the VST3
plugin's window and the LaTeX class use the same palette.
